% Part: lambda-calculus
% Chapter: introduction
% Section: composition

\documentclass[../../../include/open-logic-section]{subfiles}

\begin{document}

\olfileid{lam}{int}{com}
\olsection{الدوال \usetoken{S}{lambda definable} مغلقة تحت التركيب}

\begin{lem}
التركيب يحفظ كون الدوال~!!{lambda definable}.
\end{lem}

\begin{proof}
لتكن~$f$ حاصلة بالتركيب من~$h$ و$g_0$، \dots،~$g_{k-1}$.
ولنفرض أن~$h$ و$g_0$، \dots،~$g_{k-1}$ هي !!{lambda defined}
بالحدود~$H$ و$G_0$، \dots،~$G_{k-1}$، بالترتيب. فالمطلوب حد~$F$
!!{lambda define}s~$f$، ويحصل المطلوب بتعريف~$F$ على الوجه الآتي:
\[
F(x_0, \dots, x_{l-1}) = H(G_0(x_0, \dots, x_{l-1}),
\dots, G_{k-1}(x_0, \dots, x_{l-1})).
\]
فهذا يبيّن ملاءمة لغة حساب لامبدا لتمثيل تركيب الدوال.
\end{proof}

\end{document}

